\section{Design e Arquitetura do Sistema}
\label{sec:design}

O design e a arquitetura de um sistema de software são os alicerces sobre os quais sua funcionalidade, manutenibilidade e escalabilidade são construídas. Para um sistema distribuído e assíncrono como a plataforma de processamento de mídias, a escolha de uma arquitetura robusta não é apenas uma boa prática, mas um requisito fundamental para o sucesso do projeto. Decisões tomadas nesta fase impactam diretamente a capacidade do sistema de evoluir, de se adaptar a novas tecnologias e de resistir a falhas.

Esta seção detalha as decisões de design e a arquitetura de alto nível que foram implementadas. Começando com uma visão arquitetural abrangente, que descreve os padrões de design e os princípios que guiaram o desenvolvimento, passando pela modelagem visual do sistema com diagramas UML, e culminando no design do esquema do banco de dados que sustenta a persistência dos dados. O objetivo é fornecer uma compreensão clara e profunda de como o sistema foi estruturado para atender aos seus requisitos funcionais e, crucialmente, aos seus atributos de qualidade não funcionais, como extensibilidade, resiliência e observabilidade.
